% GENERATED FILE. Edit canonical lessons, then run npm run generate:books.
% canonical-source-hash: fnv1a64:a6283c3d4ed4f4a7
% canonical-lessons: TE-C442-patta, TE-C442-metro, TE-C442-rakapokalu, TE-C442-kikkirisipovu, TE-C442-pattalu-tappu

\chapter{On the Rails}
\label{ch:te-on-the-rails}
\hlchaptermodality{442}

\begin{chapteropening}
\textbf{By the end of this chapter:} \emph{I can say a rail, the metro, traffic, to be packed and to derail.}

Complete the last lesson of chapter 442: I can say a rail, the metro, traffic, to be packed and to derail.
\end{chapteropening}

\section[\texorpdfstring{\te{పట్టా} (paṭṭā)}{పట్టా (paṭṭā)}]{\te{పట్టా} (paṭṭā) — a rail; a track}

\label{lesson:TE-C442-patta}



\begin{quote}
Before the new one: say the Telugu for security, then the Telugu for a certificate.
\end{quote}

\subsection*{You'll want to know: \te{పట్టా}}
\textbf{\te{పట్టా}} — \emph{paṭṭā} — \textquotedblleft{}a rail; a track\textquotedblright{}.

\begin{quote}
\te{పిల్లలు} \te{రైలు} \te{పట్టాల} \te{దగ్గర} \te{ఆడకూడదు}.
\end{quote}

\begin{quote}
\emph{pillalu railu paṭṭāla daggara āḍakūḍadu.}
\end{quote}

\begin{quote}
Children must not play near the railway tracks.
\end{quote}

Plural: \te{పట్టాలు}, \emph{paṭṭālu}.

The rails of a line are \te{పట్టాలు}. \te{పట్టా} is also a strip of land and a title deed.

\begin{cousinweb}
\textbf{\te{పట్టా}} (\emph{paṭṭā}) is usually traced to Sanskrit \emph{paṭṭa}, \textquotedblleft{}a slab, a tablet, a strip\textquotedblright{} (CDIAL 7699); a rail is a long strip of iron.
\end{cousinweb}

\begin{grammarlens}[title={What you've built}]
One more word for getting around.
\end{grammarlens}

\subsection*{Your turn}
\begin{itemize}
\raggedright
  \item \emph{Say it:} \emph{paṭṭā}
  \item \emph{Say it:} \emph{paṭṭā}, once more
  \item \emph{Say it:} \emph{pillalu railu paṭṭāla daggara āḍakūḍadu.} — Children must not play near the railway tracks.
  \item \emph{Recall:} say the Telugu for security, then the Telugu for a certificate, then say \emph{paṭṭā} again
\end{itemize}

\begin{tcolorbox}[breakable,colback=teal!4,colframe=teal!35!black,title={Before you move on}]
What does \te{పట్టా} mean? (\textquotedblleft{}A rail; a track\textquotedblright{}.) Say it once more.
\end{tcolorbox}
\section[\texorpdfstring{\te{మెట్రో} (meṭrō)}{మెట్రో (meṭrō)}]{\te{మెట్రో} (meṭrō) — the metro}

\label{lesson:TE-C442-metro}



\begin{quote}
Before the new one: say the Telugu for a certificate, then the Telugu for a rail.
\end{quote}

\subsection*{You'll want to know: \te{మెట్రో}}
\textbf{\te{మెట్రో}} — \emph{meṭrō} — \textquotedblleft{}the metro\textquotedblright{}.

\begin{quote}
\te{నగరంలో} \te{మెట్రో} \te{రైలు} \te{చాలా} \te{వేగంగా} \te{వెళ్తుంది}.
\end{quote}

\begin{quote}
\emph{nagaraṁlō meṭrō railu cālā vēgaṅgā veḷtundi.}
\end{quote}

\begin{quote}
The metro train in the city goes very fast.
\end{quote}

People say \te{మెట్రో} or \te{మెట్రో} \te{రైలు}; a station on it is a \te{మెట్రో} \te{స్టేషన్}.

\begin{cousinweb}
\textbf{\te{మెట్రో}} (\emph{meṭrō}) is a loan from English \emph{metro}, short for a city railway.
\end{cousinweb}

\begin{grammarlens}[title={What you've built}]
One more word for getting around.
\end{grammarlens}

\subsection*{Your turn}
\begin{itemize}
\raggedright
  \item \emph{Say it:} \emph{meṭrō}
  \item \emph{Say it:} \emph{meṭrō}, once more
  \item \emph{Say it:} \emph{nagaraṁlō meṭrō railu cālā vēgaṅgā veḷtundi.} — The metro train in the city goes very fast.
  \item \emph{Recall:} say the Telugu for a certificate, then the Telugu for a rail, then say \emph{meṭrō} again
\end{itemize}

\begin{tcolorbox}[breakable,colback=teal!4,colframe=teal!35!black,title={Before you move on}]
What does \te{మెట్రో} mean? (\textquotedblleft{}The metro\textquotedblright{}.) Say it once more.
\end{tcolorbox}
\section[\texorpdfstring{\te{రాకపోకలు} (rākapōkalu)}{రాకపోకలు (rākapōkalu)}]{\te{రాకపోకలు} (rākapōkalu) — comings and goings; traffic}

\label{lesson:TE-C442-rakapokalu}



\begin{quote}
Before the new one: say the Telugu for a rail, then the Telugu for the metro.
\end{quote}

\subsection*{You'll want to know: \te{రాకపోకలు}}
\textbf{\te{రాకపోకలు}} — \emph{rākapōkalu} — \textquotedblleft{}comings and goings; traffic\textquotedblright{}.

\begin{quote}
\te{వరద} \te{వల్ల} \te{ఆ} \te{రోడ్డు} \te{మీద} \te{రాకపోకలు} \te{నిలిచాయి}.
\end{quote}

\begin{quote}
\emph{varada valla ā rōḍḍu mīda rākapōkalu nilicāyi.}
\end{quote}

\begin{quote}
Because of the flood, traffic on that road stopped.
\end{quote}

\te{రాక}, coming, + \te{పోక}, going, + the plural -\te{లు}. The news uses it for the traffic on a road.

\begin{cousinweb}
\textbf{\te{రాకపోకలు}} (\emph{rākapōkalu}) is native Dravidian: \te{రాక}, \textquotedblleft{}coming\textquotedblright{}, the noun of \te{రా}, \textquotedblleft{}to come\textquotedblright{} (DEDR 5270), and \te{పోక}, \textquotedblleft{}going\textquotedblright{}, the noun of \te{పోవు}, \textquotedblleft{}to go\textquotedblright{} (DEDR 4572), with the plural -\te{లు}.
\end{cousinweb}

\begin{grammarlens}[title={What you've built}]
One more word for getting around.
\end{grammarlens}

\subsection*{Your turn}
\begin{itemize}
\raggedright
  \item \emph{Say it:} \emph{rākapōkalu}
  \item \emph{Say it:} \emph{rākapōkalu}, once more
  \item \emph{Say it:} \emph{varada valla ā rōḍḍu mīda rākapōkalu nilicāyi.} — Because of the flood, traffic on that road stopped.
  \item \emph{Recall:} say the Telugu for a rail, then the Telugu for the metro, then say \emph{rākapōkalu} again
\end{itemize}

\begin{tcolorbox}[breakable,colback=teal!4,colframe=teal!35!black,title={Before you move on}]
What does \te{రాకపోకలు} mean? (\textquotedblleft{}Comings and goings; traffic\textquotedblright{}.) Say it once more.
\end{tcolorbox}
\section[\texorpdfstring{\te{కిక్కిరిసిపోవు} (kikkirisipōvu)}{కిక్కిరిసిపోవు (kikkirisipōvu)}]{\te{కిక్కిరిసిపోవు} (kikkirisipōvu) — to be packed; to be crammed}

\label{lesson:TE-C442-kikkirisipovu}



\begin{quote}
Before the new one: say the Telugu for the metro, then the Telugu for comings and goings.
\end{quote}

\subsection*{You'll want to know: \te{కిక్కిరిసిపోవు}}
\textbf{\te{కిక్కిరిసిపోవు}} — \emph{kikkirisipōvu} — \textquotedblleft{}to be packed; to be crammed\textquotedblright{}.

\begin{quote}
\te{సాయంత్రం} \te{బస్సు} \te{జనంతో} \te{కిక్కిరిసిపోయింది}.
\end{quote}

\begin{quote}
\emph{sāyantraṁ bassu janantō kikkirisipōyindi.}
\end{quote}

\begin{quote}
The evening bus was packed with people.
\end{quote}

Past (it …): \te{కిక్కిరిసిపోయింది}, \emph{kikkirisipōyindi}.

What fills the place takes -\te{తో}: \te{జనంతో}, with people. \te{పోవు} at the end marks a change that has happened completely. Its subject is a thing, so the past shown is the third person: \te{కిక్కిరిసిపోయింది}, it was packed.

\begin{cousinweb}
\textbf{\te{కిక్కిరిసిపోవు}} (\emph{kikkirisipōvu}) is a Telugu compound whose first part, \te{కిక్కిరియు}, \textquotedblleft{}to be crammed\textquotedblright{}, has an origin that is not settled; the second part is the native \te{పోవు}, \textquotedblleft{}to go\textquotedblright{} (DEDR 4572), which marks a change that is complete.
\end{cousinweb}

\begin{grammarlens}[title={What you've built}]
One more word for getting around.
\end{grammarlens}

\subsection*{Your turn}
\begin{itemize}
\raggedright
  \item \emph{Say it:} \emph{kikkirisipōvu}
  \item \emph{Say it:} \emph{kikkirisipōvu}, once more
  \item \emph{Say it:} \emph{sāyantraṁ bassu janantō kikkirisipōyindi.} — The evening bus was packed with people.
  \item \emph{Recall:} say the Telugu for the metro, then the Telugu for comings and goings, then say \emph{kikkirisipōvu} again
\end{itemize}

\begin{tcolorbox}[breakable,colback=teal!4,colframe=teal!35!black,title={Before you move on}]
What does \te{కిక్కిరిసిపోవు} mean? (\textquotedblleft{}To be packed; to be crammed\textquotedblright{}.) Say it once more.
\end{tcolorbox}
\section[\texorpdfstring{\te{పట్టాలు} \te{తప్పు} (paṭṭālu tappu)}{పట్టాలు తప్పు (paṭṭālu tappu)}]{\te{పట్టాలు} \te{తప్పు} (paṭṭālu tappu) — to derail}

\label{lesson:TE-C442-pattalu-tappu}



\begin{quote}
Before the new one: say the Telugu for comings and goings, then the Telugu for to be packed.
\end{quote}

\subsection*{You'll want to know: \te{పట్టాలు} \te{తప్పు}}
\textbf{\te{పట్టాలు} \te{తప్పు}} — \emph{paṭṭālu tappu} — \textquotedblleft{}to derail\textquotedblright{}.

\begin{quote}
\te{నిన్న} \te{రాత్రి} \te{ఒక} \te{సరుకు} \te{రైలు} \te{పట్టాలు} \te{తప్పింది}.
\end{quote}

\begin{quote}
\emph{ninna rātri oka saruku railu paṭṭālu tappindi.}
\end{quote}

\begin{quote}
Last night a goods train went off the rails.
\end{quote}

Past (it …): \te{పట్టాలు} \te{తప్పింది}, \emph{paṭṭālu tappindi}.

Word for word: to miss the rails. It is also said of a plan that goes wrong. Its subject is a thing, so the past shown is the third person: \te{పట్టాలు} \te{తప్పింది}, it went off the rails.

\begin{cousinweb}
\textbf{\te{పట్టాలు} \te{తప్పు}} (\emph{paṭṭālu tappu}) is a Telugu phrase built from \te{పట్టాలు}, \textquotedblleft{}rails\textquotedblright{}, usually traced to Sanskrit \emph{paṭṭa} (CDIAL 7699), and the native verb \te{తప్పు}, \textquotedblleft{}to slip, to miss, to go wrong\textquotedblright{} (DEDR 3071).
\end{cousinweb}

\begin{grammarlens}[title={What you've built}]
One more word for getting around.
\end{grammarlens}

\subsection*{Your turn}
\begin{itemize}
\raggedright
  \item \emph{Say it:} \emph{paṭṭālu tappu}
  \item \emph{Say it:} \emph{paṭṭālu tappu}, once more
  \item \emph{Say it:} \emph{ninna rātri oka saruku railu paṭṭālu tappindi.} — Last night a goods train went off the rails.
  \item \emph{Recall:} say the Telugu for comings and goings, then the Telugu for to be packed, then say \emph{paṭṭālu tappu} again
\end{itemize}

\begin{tcolorbox}[breakable,colback=teal!4,colframe=teal!35!black,title={Before you move on}]
What does \te{పట్టాలు} \te{తప్పు} mean? (\textquotedblleft{}To derail\textquotedblright{}.) Say it once more.
\end{tcolorbox}
